\section{Branching and neighbouring factor dimensions}\label{sec:mixed-branching}

We retain the mixed-blowup construction $Y(G,b)$ of
Section~\ref{sec:mixed}, with distinct original rays assigned to the
incidences at each vertex. Write $d_v=\deg_G(v)$. Throughout this
section assume $b_v\ge d_v$ and $b_v\ge1$ at every vertex, so that
$Y(G,b)$ is Fano. A vertex $v$ is
\emph{saturated} if $b_v=d_v$; equivalently, exactly one of its original
rays is unassigned. The degree of its unassigned divisor depends on
the geometry beyond its incident centres. We first determine stability
for saturated stars, including the stable degree-four and degree-five
examples. We then show that replacing the rest of the geometry by
independent leaves of the same neighbouring dimensions can only decrease
this degree. This comparison allows cycles and does not require
saturation at the other vertices.

\subsection{The exact transition for stars}

Let $Y_m$ be the mixed blowup associated with a star with $m\ge1$ edges,
central dimension $m$, and dimension one at each leaf. In the product
lattice, choose bases $e_1,\ldots,e_m$ for the central factor and $u_i$
for leaf $i$. Put
\[
 e_0=-\sum_{i=1}^m e_i,\qquad s_i=e_i+u_i.
\]
Thus $Y_m$ is obtained from $\PP^m\times(\PP^1)^m$ by subdividing the
$m$ original cones $\langle e_i,u_i\rangle$. Let $A_m$ denote the
normalized divisor degree of $e_0$, with the normalization used in
Section~\ref{sec:mixed}.

\begin{theorem}\label{branch:stars}
For every $m\ge1$, the variety $Y_m$ is smooth Fano, has dimension $2m$
and Picard number $2m+1$, and admits no torus-equivariantly locally
trivial bundle presentation with positive-dimensional complete fibre
and positive-dimensional base. Its tangent bundle is anticanonically
stable exactly when $1\le m\le5$. For every $m\ge6$, the line
$\CC e_0$ defines a strictly destabilizing saturated subsheaf.
\end{theorem}

The proof, including the finite sign calculation and the uniform tail
estimate, is given in Appendix~\ref{sec:blowup-comparisons}.


For the two stable stars with higher valence, the central degrees and
total exceptional contributions are
\[
 A_4=\frac{44080}{47119},\qquad
 \sum_iJ_{0,i}=\frac{31924}{47119}<1;
 \qquad
 A_5=\frac{100780}{103289},\qquad
 \sum_iJ_{0,i}=\frac{88235}{103289}<1.
\]
These values follow from the table and \eqref{mixed:flux-identity}.
They show that neither four nor five incident blowups alone forces the
radial foliation of Proposition~\ref{intrinsic:foliation} to destabilize.
The neighbouring slice profiles, rather than the number of contributions,
determine whether their sum reaches one.

The same line is the first Harder--Narasimhan term for $m\ge6$
(Corollary~\ref{branch:hn}).

\subsection{Comparison at a saturated vertex}

The star calculation gives a necessary condition for stability in the
entire simple-graph construction. In the next statement, the graph need
not be connected and may contain cycles.

\begin{theorem}\label{branch:comparison}
Let $G$ be a finite simple graph with positive dimensions $b_w\ge d_w$
at every vertex, and use injective incidence assignments in the
construction of $Y(G,b)$. Suppose $v$ is saturated and
$b_v=d_v=m\ge1$. Let $A_v$ be the normalized divisor degree of its unique
unassigned original ray. Enumerate the neighbors as $w_1,\ldots,w_m$
and put $c_i=b_{w_i}$. Define
\[
 \phi_c(a)=\frac{(c+\min(a,1))^c}{c!},\qquad
 Z_{\mathbf c}(t)=\int_{\substack{a_i\ge0\\\sum_i a_i\le t}}
                  \prod_{i=1}^m\phi_{c_i}(a_i)\,da.
\]
Then
\[
 A_v\ge\frac{Z'_{\mathbf c}(m+1)}{Z_{\mathbf c}(m+1)}
       \ge A_m.
\]
The first ratio is nondecreasing in each neighboring dimension $c_i$.
The dimensions at vertices other than $v$ may exceed their degrees.
\end{theorem}

\begin{proof}
At every vertex $w$, choose an unassigned original ray as the negative
sum ray, and use the other $b_w$ rays as a lattice basis. This is
possible because $b_w\ge d_w$. In affine coordinates
$a_{w,j}=1+\langle x,r_{w,j}\rangle$, the anticanonical polytope is
\begin{equation}\label{branch:graph-polytope}
 a_{w,j}\ge0,\qquad
 \sum_{j=1}^{b_w}a_{w,j}\le b_w+1,\qquad
 a_{w,e}+a_{z,e}\ge1\quad(e=\{w,z\}).
\end{equation}
Here $a_{w,e}$ is the coordinate assigned to the incidence $(w,e)$.
The coordinate change preserves lattice volume.

Write $a=(a_1,\ldots,a_m)$ for the coordinates at $v$. Before imposing
the simplex inequality at $v$, integrate over all other vertices and
denote the resulting volume by $W(a)$, for every $a_i\ge0$.
It is bounded by the product of the other simplex volumes and is
coordinatewise nondecreasing, since increasing $a_i$ only relaxes one
inequality. It is constant in each coordinate once that coordinate is
at least one.

Moreover, $W(0)>0$. At every other vertex choose its assigned
coordinates equal to $1+\varepsilon$ and its remaining coordinates
equal to $\varepsilon$. For sufficiently small $\varepsilon>0$, the
sum at $w$ is $d_w+b_w\varepsilon<b_w+1$, and all mixed inequalities
are strict. An open set of such choices has positive volume.
Monotonicity gives $W(a)>0$ throughout the nonnegative orthant.
The function $W$ is continuous: its domain is contained in a fixed
compact product of simplexes, and under a convergent change of $a$
the defining indicators converge away from finitely many affine
hyperplanes of measure zero.

We first prove that
\begin{equation}\label{branch:ratio}
 R(a)=\frac{W(a)}{\prod_{i=1}^m\phi_{c_i}(a_i)}
\end{equation}
is coordinatewise nondecreasing. Fix every central coordinate except
$a_i$, and let $w$ be the other endpoint of its edge. Put $b=b_w$.
At $w$, let $x$ be the coordinate assigned to this edge and let
$z=(z_1,\ldots,z_{b-1})$ be its remaining coordinates. Injectivity of
the incidence assignments ensures that no other mixed inequality
uses $x$. Simplicity ensures that $w$ has no second edge to $v$.

Temporarily omit the simplex inequality at $w$ and integrate over
every vertex other than $v,w$. The resulting function $\Psi(z)$ is
nonnegative and coordinatewise nondecreasing on the nonnegative
orthant: increasing a component of $z$ only enlarges the feasible set.
It may depend on the other fixed central coordinates, but it is
independent of $x$ and $a_i$. Edges among the other vertices remain
inside this integral, so no independence of neighbouring components
is required.

For $0\le a_i\le1$, the interval for $x$ is
$[1-a_i,b+1-\sum_jz_j]$. With $s=b+a_i$, integration over this interval
and then the substitution $z=sy$ give
\begin{align}
 W(a)
 &=\int_{\substack{z_j\ge0\\\sum_jz_j\le s}}
       (s-\textstyle\sum_jz_j)\Psi(z)\,dz \notag\\
 &=s^b\int_{\substack{y_j\ge0\\\sum_jy_j\le1}}
       (1-\textstyle\sum_jy_j)\Psi(sy)\,dy.
 \label{branch:neighbour-integral}
\end{align}
When $b=1$, the last integral is zero-dimensional and the same identity
holds. The integral factor is nondecreasing in $s$. Thus, with the
other coordinates fixed and $0\le a\le a'\le1$,
\[
 \frac{W(a')}{W(a)}
 \ge\left(\frac{b+a'}{b+a}\right)^b
 =\frac{\phi_b(a')}{\phi_b(a)}.
\]
This proves monotonicity
of $R$ below one in each coordinate. Above one both $W$ and $\phi_{c_i}$
are constant, and continuity gives the assertion across one.
The function $R$ is also bounded, since each $\phi_{c_i}$ has a positive lower bound.

To compare the simplex integrals, introduce independent random variables
$X_1,\ldots,X_m$ with respective densities
\[
 f_i(x)=C_i e^{-x}\phi_{c_i}(x)\mathbf 1_{\{x\ge0\}},
\]
where $C_i>0$ is the normalizing constant. Each density is log-concave.
On $(0,1)$ its logarithmic derivative is $-1+c_i/(c_i+x)$, which decreases;
at one this derivative jumps downward to $-1$ and stays constant
thereafter. Its support is the interval $[0,\infty)$.

Efron's conditional-sum theorem states that, for independent
log-concave random variables, the conditional expectation of a bounded
coordinatewise nondecreasing function is nondecreasing in their sum;
the version for general log-concave densities and any finite number
of variables is given in
\cite[Theorems~3.2--3.3]{saumardwellner2018}.
Extend $R$ to $\RR^m$ by replacing each argument by its nonnegative
part. The extension is bounded and coordinatewise nondecreasing.
The theorem therefore implies that
\[
 h(s)=\mathbb E\!\left[
 R(X_1,\ldots,X_m)\,\middle|\,X_1+\cdots+X_m=s
 \right]
\]
has a nondecreasing version. When $m=1$, the same assertion follows
directly from $h(s)=R(s)$.

On the slice $\sum_i a_i=s$, the product exponential factor equals
$e^{-s}$ and cancels from the conditional expectation. Thus $h(s)$
is the ratio of the $W$-weighted slice volume to the
$\prod_i\phi_{c_i}(a_i)$-weighted slice volume. These integrals are
continuous for $s>0$, by rescaling the compact sum simplex to a fixed
simplex and applying dominated convergence. The denominator is
positive. Their ratio consequently agrees everywhere on $(0,\infty)$
with the nondecreasing version supplied by the theorem.

Set
\[
 F(t)=\int_{\substack{a_i\ge0\\\sum_i a_i\le t}}W(a)\,da,
 \qquad q(t)=Z'_{\mathbf c}(t)>0.
\]
Using the lattice slice measure obtained by eliminating one coordinate,
we have
\[
 F'(t)=q(t)h(t),\qquad
 F(t)=\int_0^tq(s)h(s)\,ds
 \le h(t)Z_{\mathbf c}(t).
\]
All quantities are positive, so
\[
 \frac{F'(t)}{F(t)}\ge\frac{Z'_{\mathbf c}(t)}{Z_{\mathbf c}(t)}.
\]
At $t=m+1$, the denominator on the left is the full anticanonical
polytope volume and its numerator is the lattice volume of the
unassigned facet at $v$. Their ratio is $A_v$, proving the first comparison.

For integers $b\ge c\ge1$ and $0<a<1$,
\[
 \frac{d}{da}\log\frac{\phi_b(a)}{\phi_c(a)}
 =\frac{a(b-c)}{(b+a)(c+a)}\ge0.
\]
The ratio is constant above one. Applying the same conditional-sum
argument to products of these ratios proves monotonicity in the
neighboring dimensions. Taking every $c_i=1$ gives
$\phi_1(a)=1+\min(a,1)$ and recovers the star ratio, proving the
second comparison.
\end{proof}

\begin{corollary}\label{branch:obstruction}
In the Fano simple-graph mixed-blowup construction, a saturated vertex of
degree at least six forces anticanonical tangent instability. In
particular, if every vertex is saturated and the tangent bundle is
stable, then the maximum graph degree is at most five. The star $Y_5$
attains degree five.
\end{corollary}

\begin{proof}
Let $v$ be saturated of degree $m\ge6$. Theorems~\ref{branch:stars}
and~\ref{branch:comparison} give $A_v\ge A_m>1$. Its unassigned ray
is the negative sum of $m$ basis vectors and is not collinear with
another original ray in that factor. Original rays at other vertices
lie in different direct summands, and exceptional rays have nonzero
components in two summands. Thus its line has normalized ray weight
exactly $A_v>1$ and defines a strict rank-one destabilizer.
The last assertion follows from Theorem~\ref{branch:stars}.
\end{proof}

\begin{corollary}\label{branch:degree-five}
If a saturated degree-five vertex has at least four neighbors of factor
dimension at least two, its unassigned divisor degree satisfies
\[
 A_v\ge\frac{6244655501}{6233518181}>1,
\]
and its unassigned ray defines a strict rank-one destabilizer.
Consequently, in a saturated simple graph with stable tangent bundle,
each degree-five vertex has at least two degree-one neighbors.
\end{corollary}

\begin{proof}
Theorem~\ref{branch:comparison} reduces the assertion to the neighboring
dimensions $(2,2,2,2,1)$. Their simplex and facet volumes at capacity six are
\[
 Z(6)=\frac{6233518181}{362880},\qquad
 Z'(6)=\frac{892093643}{51840}.
\]
To specify these rational calculations explicitly, let $c_{j,k}$ be
the coefficient of $q^ju^k$ in
\[
 (2u^2+2u+1-q(3u+1))^4(u+1-q).
\]
The Laplace transforms of $\phi_1,\phi_2$ are respectively
$(s+1-e^{-s})/s^2$ and
$(2s^2+2s+1-e^{-s}(3s+1))/s^3$. Thus
\[
 Z(6)=\sum_{j,k}c_{j,k}\frac{(6-j)^{14-k}}{(14-k)!},\qquad
 Z'(6)=\sum_{j,k}c_{j,k}\frac{(6-j)^{13-k}}{(13-k)!},
\]
where $0\le j\le5$ and $0\le k\le9$. These finite sums give the displayed
values. The unassigned ray is the only ray on its line, by the same
direct-summand argument as in Corollary~\ref{branch:obstruction}.
Its weight exceeds one, so the line strictly destabilizes.
In the saturated case neighboring dimensions equal neighboring degrees;
at most three of the five neighbors can therefore have degree at least two.
\end{proof}

\begin{corollary}\label{branch:subdivision}
A saturated five-arm tree is unstable whenever at least four of its arms
have length at least two. In particular, subdivision of four arms of the
stable five-arm star produces an unstable mixed blowup, and further
subdivision of those arms preserves this obstruction.
\end{corollary}

\begin{proof}
The central vertex has degree five and at least four degree-two neighbors,
so Corollary~\ref{branch:degree-five} applies.
\end{proof}

Here subdivision refers to the indexing graph: its factor dimensions
and the dimension of the variety change. It is not a fan subdivision in
fixed dimension. The obstruction rules out a general stability theorem
obtained merely by making every arm of a mixed-blowup tree sufficiently
long. It concerns saturated vertices and places no corresponding degree
bound on vertices whose dimensions exceed their degrees.
